\documentclass[12pt]{article}
\usepackage{amsmath, amssymb, enumitem, fancyhdr, graphicx, libertine, nth,
  pgfplots, tikz, xcolor}
\pgfplotsset{compat=1.11}
\usepackage[margin=1in]{geometry}
\usepackage[libertine]{newtxmath}
\pagestyle{fancy}
\chead{\emph{EC201 -- Problem Set 1}}

% Define new topic command
\newcommand{\topic}[1]{\medskip\begin{center}\large\textsc{#1}\normalsize
  \end{center}}

% Define question counter:
\newcounter{question}[section]
\newcommand{\question}[1]{\refstepcounter{question}
  \subsubsection*{Question~\thequestion~-- #1}}

% Subquestions are enumerate with letters
\newenvironment{subquestions}{
  \begin{enumerate}[label=(\roman*)]}{\end{enumerate}}

\begin{document}

\topic{Mathematics Review}
\question{Differentiation}
Differentiate the following functions, i.e.~find $\frac{df\left( x
\right)}{dx}$:
\begin{subquestions}
  \item $f\left( x \right)= 4x^2$
  \item $f\left( x \right)= 201$
  \item $f\left( x \right)= 2x + 4x^3$
\end{subquestions}

\question{Partial Differentiation}
Find the partial derivatives of the following functions, i.e.~find
$\frac{\partial f\left( x_1,x_2 \right)}{\partial x_1}$ and
$\frac{\partial f\left( x_1,x_2 \right)}{\partial x_2}$:
\begin{subquestions}
  \item $f\left( x_1,x_2 \right)=x_1^{\frac{1}{3}}x_2^{\frac{2}{3}}$
  \item $f\left( x_1,x_2 \right)=2x_1 + 3x_2^2$
  % \item $f\left( x_1,x_2 \right)=\sqrt{x_1x_2}\;\;\;$ (Hint:
  %   $\sqrt{x}$ is the same as $x^{\frac{1}{2}}$).
\end{subquestions}

\question{Optimization}
For each of the functions below, do the following:
\begin{itemize}
  \item Sketch the graph of the function.
  \item Find $\frac{df\left( x
\right)}{dx}$
  \item Find the extreme point, i.e.~at what $x^{\star}$ is $\frac{df\left(
    x^{\star} \right)}{dx}=0$?
  \item Is the extreme point a maximum or a minimum? (To do this you will need
    to find the second derivative.)
\end{itemize}
\begin{subquestions}
  \item $f\left( x \right) = 8x - 2x^2$
  \item $f\left( x \right) = x^2$
\end{subquestions}

\topic{The Budget Constraint}
\question{Two Price Changes}
An individual's budget line is given by:
\[
  p_1x_1 + p_2x_2 = m
\]
Suppose the price of good 1 doubled and the price of good 2 halved.
Sketch the original and new budget constraints.

\question{Price and Income Changes}
Suppose next year the prices of goods 1 and 2 will
increase by 5\% and your salary will also increase by 5\%. Describe how this
affects your budget line.

\topic{Preferences}
\question{Indifference Curves}
Suppose you enjoy consuming good 1 but you are not at all interested in good 2 (you are neutral about good 2). If you ever get any of good 2 you just give it away.  What would your indifference curves look like? Sketch them.

\topic{Utility}
\question{Indifference Curves and the Marginal Rate of Substitution}
Answer the following questions about these utility functions:
\begin{subquestions}
  \item $u\left(x_1, x_2  \right)=x_1 + 2x_2$.
    \begin{itemize}
      \item Sketch the indifference curves of the utility function for
        fixed levels of utility $k=1$, $k=2$ and $k=3$.
      \item Calculate the marginal rate of substitution ($MRS$) between the two
        goods, where: \[ MRS = -\frac{\frac{\partial u\left( x_1,x_2
        \right)}{\partial x_1}}{\frac{\partial u\left( x_1,x_2
      \right)}{\partial x_2}} \]
      \item Can you come up with an example of two goods where this
      utility function would be reasonable?
    \end{itemize}
  \item $u\left( x_1,x_2 \right)=\min\left\{\frac{1}{2}x_1, x_2\right\}$
    \begin{itemize}
      \item Sketch the indifference curves of the utility function for
        fixed levels of utility $k=1$, $k=2$ and $k=3$.
      \item What is the marginal rate of substitution when $x_1 = 3$ and
        $x_2=1$? Interpret this number.
    \end{itemize}
  \item $u\left( x_1, x_2 \right)=x_1^{\frac{1}{3}}x_2^{\frac{2}{3}}$
    \begin{itemize}
      \item Sketch the indifference curves of the utility function for
        fixed levels of utility $k=1$, $k=2$ and $k=3$.
      \item Calculate the marginal rate of substitution between the two
        goods.
    \end{itemize}
\end{subquestions}
\end{document}
